\documentclass[12pt, twoside]{article}
\input{../preamble}

\fancyhead[LE]{\thepage}
\fancyhead[RO]{Markscheme}
\fancyhead[LO]{La Scuola d'Italia / Dr.\ Huson / IB Math \\* 24 April 2026}

\begin{document}
\subsubsection*{5.17 Classwork: Logarithms Practice --- Markscheme}

\begin{enumerate}[itemsep=0.3cm]

%--- Problem 1: Evaluate basic logarithms [12 marks] ---
%--- Source: adapted from 7-1CW_Intro_logs Q1 ---
\item[1.] [Maximum mark: 12]

Evaluate each of the following.

\textbf{Solution:}

\textbf{(a)} $\dfrac{1}{64} = 4^{-3}$ \hfill \textbf{M1}

$\log_4 \dfrac{1}{64} = -3$ \hfill \textbf{A1}

\medskip

\textbf{(b)} $5 = 25^{1/2}$ \hfill \textbf{M1}

$\log_{25} 5 = \dfrac{1}{2}$ \hfill \textbf{A1}

\medskip

\textbf{(c)} $\sqrt{32} = 2^{5/2}$ \hfill \textbf{M1}

$\log_2 \sqrt{32} = \dfrac{5}{2}$ \hfill \textbf{A1}

\medskip

\textbf{(d)} $125 = 5^3$ \hfill \textbf{M1}

$\log_5 125 = 3$ \hfill \textbf{A1}

\medskip

\textbf{(e)} $\dfrac{1}{4} = 16^{-1/2}$ \hfill \textbf{M1}

$\log_{16} \dfrac{1}{4} = -\dfrac{1}{2}$ \hfill \textbf{A1}

\medskip

\textbf{(f)} $\sqrt{81} = 9 = 3^2$ \hfill \textbf{M1}

$\log_3 \sqrt{81} = 2$ \hfill \textbf{A1}

\newpage

%--- Problem 2: Write logarithmic expressions in terms of a parameter [12 marks] ---
%--- Source: adapted from 7-1CW_Intro_logs Q2 ---
\item[2.] [Maximum mark: 12]

\textbf{(a)} Let $c = \log_3 m$, where $m > 0$. Write each expression
in terms of $c$. \hfill [6]

\textbf{Solution:}

\textbf{(i)} $\log_3 m^2 = 2\log_3 m$ \hfill \textbf{M1}

$= 2c$ \hfill \textbf{A1}

\medskip

\textbf{(ii)} $\log_3(27m) = \log_3 27 + \log_3 m$ \hfill \textbf{M1}

$= 3 + c$ \hfill \textbf{A1}

\medskip

\textbf{(iii)} $\log_9 m = \dfrac{\log_3 m}{\log_3 9}$ \hfill \textbf{M1}

$= \dfrac{c}{2}$ \hfill \textbf{A1}

\medskip

\textbf{(b)} Let $a = \log_2 m$, where $m > 0$. Write each expression
in terms of $a$. \hfill [6]

\textbf{Solution:}

\textbf{(i)} $\log_2 m^3 = 3\log_2 m$ \hfill \textbf{M1}

$= 3a$ \hfill \textbf{A1}

\medskip

\textbf{(ii)} $\log_2(8m) = \log_2 8 + \log_2 m$ \hfill \textbf{M1}

$= 3 + a$ \hfill \textbf{A1}

\medskip

\textbf{(iii)} $\log_4 m = \dfrac{\log_2 m}{\log_2 4}$ \hfill \textbf{M1}

$= \dfrac{a}{2}$ \hfill \textbf{A1}

\newpage

%--- Problem 3: Expand and condense logarithmic expressions [12 marks] ---
%--- Source: adapted from 7-1CW_Intro_logs Q3 ---
\item[3.] [Maximum mark: 12]

Use the laws of logarithms.

\textbf{(a)} Express each in terms of $\log x$ and $\log y$. \hfill [6]

\textbf{Solution:}

\textbf{(i)} $\log(100x) = \log 100 + \log x$ \hfill \textbf{M1}

$= 2 + \log x$ \hfill \textbf{A1}

\medskip

\textbf{(ii)}
$\log \!\left(\dfrac{x^2}{y}\right) = \log x^2 - \log y$ \hfill \textbf{M1}

$= 2\log x - \log y$ \hfill \textbf{A1}

\medskip

\textbf{(iii)}
$\log \sqrt{xy^3} = \dfrac{1}{2}\log(xy^3)$ \hfill \textbf{M1}

$= \dfrac{1}{2}\log x + \dfrac{3}{2}\log y$ \hfill \textbf{A1}

\medskip

\textbf{(b)} Write each expression as a single logarithm. \hfill [6]

\textbf{Solution:}

\textbf{(i)} $\log_2 x + \log_2 y = \log_2(xy)$ \hfill \textbf{M1A1}

\medskip

\textbf{(ii)}
$2\log_7 a - \log_7 b = \log_7 a^2 - \log_7 b$ \hfill \textbf{M1}

$= \log_7\!\left(\dfrac{a^2}{b}\right)$ \hfill \textbf{A1}

\medskip

\textbf{(iii)}
$\dfrac{1}{2}\ln p + \ln q = \ln p^{1/2} + \ln q$ \hfill \textbf{M1}

$= \ln(q\sqrt{p})$ \hfill \textbf{A1}

\emph{Accept $\ln(\sqrt{p}\,q)$ or $\ln(p^{1/2}q)$.}

\newpage

%--- Problem 4: Solve basic logarithmic equations [12 marks] ---
%--- Source: adapted from 7-1CW_Intro_logs Q4 ---
\item[4.] [Maximum mark: 12]

Solve each equation. Give exact answers.

\textbf{Solution:}

\textbf{(a)} $\log_4(x + 9) = \dfrac{1}{2}\log_2(x + 9)$ \hfill \textbf{M1}

$2\log_2(x - 3) = \log_2(x + 9)$

$(x - 3)^2 = x + 9$ \hfill \textbf{M1}

$x^2 - 7x = 0$

$x = 0$ or $x = 7$

Domain requires $x > 3$, so $x = 7$ \hfill \textbf{A1}

\medskip

\textbf{(b)} $\log_3 x = 2$

$x = 3^2$ \hfill \textbf{M1}

$x = 9$ \hfill \textbf{A1A1}

\medskip

\textbf{(c)} $\log_5(2x) = 1$

$2x = 5^1$ \hfill \textbf{M1}

$x = \dfrac{5}{2}$ \hfill \textbf{A1A1}

\medskip

\textbf{(d)}
$\log_2(x + 1) + \log_2(x - 1) = \log_2((x + 1)(x - 1))$ \hfill \textbf{M1}

$x^2 - 1 = 8$ \hfill \textbf{M1}

$x^2 = 9$

$x = \pm 3$ 

\emph{Reject $x = -3$ since $x > 1$.}

$x = 3$ \hfill \textbf{A1}


\end{enumerate}
\end{document}
